\section{Capítulo 6: Conclusiones}
\label{conclusiones}

\subsection{Conclusiones}

El trabajo se propuso integrar una herramienta basada en un modelo de lenguaje de gran tamaño a una arquitectura Unified Namespace para la analítica de datos de un caso de Industria 4.0, y esa integración se logró. En el entorno de prueba, un usuario pregunta en español y con el vocabulario de la planta por el estado actual de una línea, por su historial o por ambos, y recibe una respuesta respaldada por los datos del UNS y de la base \texttt{mes\_core}. En las 75 ejecuciones de la validación, 71 fueron correctas (94,7~\%) y ninguna contuvo una cifra, un nombre o una causa inventados. La validación comprueba que la integración funciona para este caso. No compara modelos ni permite decir que otro modelo se comportaría igual.

Cada objetivo específico aportó una pieza. El Capítulo~\ref{cap:caso_estudio} caracterizó la planta de empaque y fijó tres requerimientos y ocho desafíos que sirvieron de criterio a todo lo posterior (OE1). El Capítulo~\ref{tecnicasdiseñouns} entregó la arquitectura UNS, con una jerarquía y unos UDT que describen cada variable, de modo que el dato llega acompañado de su significado (OE2). El Capítulo~\ref{seleccionllm} comparó los mecanismos de integración frente a las prioridades del caso y eligió MCP por su capacidad de exponer una jerarquía compleja y datos en tiempo real, aceptando que su control de acceso era menos maduro que el de REST y dejando ese riesgo para la implementación (OE3). El Capítulo~\ref{implementacionintegracion} lo resolvió con decisiones de diseño: dos servidores MCP con transporte STDIO, Tools de solo lectura que ejecutan consultas parametrizadas sobre una lista cerrada de tablas, sin SQL libre, y Resources que entregan al modelo la jerarquía y el glosario de la planta (OE4). El Capítulo~\ref{validacionsistema} contrastó el sistema con 25 preguntas cuyas respuestas de referencia se calcularon por una vía independiente de las Tools, y con los desafíos del caso (OE5). Siete de los ocho quedaron cumplidos, dos de ellos con salvedades acotadas y uno por diseño. El octavo, la adaptación del reporte al perfil del usuario (\desafio{D2.3}), no se implementó.

\subsection{Aporte del trabajo}

El aporte es un mecanismo concreto para que un LLM consulte un UNS sin sustituirlo como fuente de verdad. El significado de los datos, que en un UNS ya está en la jerarquía, los nombres y las unidades, llega al modelo a través de MCP y no hay que reconstruirlo en cada pregunta. Con ello se atiende la limitación planteada en la justificación del trabajo: centralizar los datos no basta para aprovecharlos, y es el contexto semántico lo que permite a un modelo de lenguaje interpretar un namespace ordenado. Las Tools acotadas y de solo lectura agregan otra propiedad. El modelo no accede libremente a la base ni al sistema de control, sino a un conjunto pequeño de operaciones cuyo alcance se conoce de antemano, y eso permite revisar el riesgo de la integración.

La validación mostró además dónde queda ese riesgo. Los errores observados no se originaron en una cifra inventada por el modelo. Provinieron de una Tool que dejaba fuera las órdenes del último día del rango y de una expresión, «producción del día», que admitía dos definiciones. El modelo confía en lo que devuelve la Tool y lo presenta con seguridad, de modo que la calidad de las respuestas depende de la calidad de las Tools y de los metadatos que las acompañan. Quien despliegue una integración similar debe tratar la verificación de las Tools con referencias independientes como parte del diseño.

\subsection{Alcance y límites}

Los resultados se leen dentro de su alcance. Los valores de los tags y los registros históricos se generaron con scripts que respetan la estructura, las unidades y las reglas de negocio del caso real, pero no vienen de una planta en operación. Por eso no se probó el comportamiento del sistema ante el ruido y los vacíos de datos reales, y el simulador no reproduce todas las relaciones entre variables de una planta, como la que existe entre la tasa estándar y la producción. La validación usó un banco de 25 preguntas, un único modelo y una única configuración, porque el objetivo era comprobar la integración. Los porcentajes describen este sistema sobre este banco y no se extienden a otros modelos ni a otras plantas.

\subsection{Trabajo futuro}

Lo más inmediato es cerrar lo que la validación dejó identificado. El filtro de fecha de la Tool de órdenes y la documentación de las unidades ya se corrigieron y se comprobaron con un banco adicional (Anexo~\ref{anexo:validacion_correcciones}); falta fijar la definición de producción diaria, que admitía dos lecturas, y repetir el banco completo para confirmar todas las correcciones. Después, el paso de mayor valor es repetir la validación con datos de una planta en operación, donde aparecerán los casos que la simulación no reproduce. Mientras tanto, conviene ajustar el simulador para que conserve las relaciones entre variables, como la de la tasa estándar con la producción. También queda por implementar la adaptación de las respuestas al perfil del usuario (\desafio{D2.3}), para que quien consulta reciba el nivel de detalle que su función requiere. La misma estructura de Tools y Resources puede extenderse a otras fuentes de la empresa, como los sistemas de mantenimiento o de calidad. Si el caso lo justifica, podría estudiarse la escritura controlada sobre el UNS, con confirmación humana y permisos propios. Esa ampliación obligaría a revisar la seguridad de la integración, que en este trabajo se apoyó en que todo acceso es de solo lectura.
